% =====================================================================================================
% Chương 4, PHẦN A (mở chương, mục 4.1–4.5) — bản viết lại (2026-10-11).
% Dùng làm phần đầu của \input{noi-dung/08-chuong4}; phần B (mục 4.6–4.11) nối ngay sau tệp này.
% Ghi chú cho người ghép: xem 08-chuong4-phanA-NOTES.md (nhãn cần phần B / chương khác, bằng chứng thiếu).
% Mỗi dòng "% EV:" trỏ tới dòng của EVIDENCE-REGISTER.md (A01…K52), EVIDENCE-REGISTER-ADDENDUM-1.md
% (N-01…N-30, L1a-01…09, mục 0.x) hoặc mã X3-xx (dữ kiện mã nguồn của người viết chương 3, đã được
% R23-FACTCHECK-CH3 kiểm; chưa có trong sổ, liệt kê trong NOTES).
% Không dùng gói mới: hình vẽ bằng TikZ thường (thư viện shapes.geometric đã nạp trong main.tex).
% =====================================================================================================

\chapter{Thực nghiệm và đánh giá}
\label{chuong:thucnghiem}

Chương này đo luật C được gì và mất gì trên mô phỏng FleetPy, so với điều phối không trần, luật hạn
chót và cách đệm lời hứa, rồi xét lời hứa theo dự báo và những chỗ dự báo chưa đủ. Các mục đi theo ba
phần của độ lệch giờ trả. Phần do điều phối ở mục~\ref{muc:c4-ketqua}--\ref{muc:c4-vaitro}; phần được hồ
sơ lịch sử giải thích, cùng cách đệm lời hứa, ở mục~\ref{muc:c4-dubao} và~\ref{muc:c4-dem}; phần còn lại
ở mục~\ref{muc:c4-tuyet}. Hai mục cuối kiểm tính hợp lệ của kết quả và bàn các hạn chế.
Bảng~\ref{bang:c4-lotrinh} nối ba câu hỏi nghiên cứu ở mục~\ref{muc:c1-cauhoi} với các mục trả lời chúng.
% EV: — (đoạn dẫn, không nêu số; thứ tự mục theo OUTLINE-V2 Chương 4)

\begin{table}[htbp]
\centering
\caption{Ba câu hỏi nghiên cứu và các mục của Chương~\ref{chuong:thucnghiem} trả lời chúng}
\label{bang:c4-lotrinh}
\small
\begin{tabular}{@{}p{6.3cm}p{2.4cm}p{6.3cm}@{}}
\toprule
\textbf{Câu hỏi} & \textbf{Mục} & \textbf{Đại lượng được đo}\\
\midrule
CH1. Phần dời do điều phối được đo thế nào, và trần trên phần đó giữ được gì, không giữ được gì?
& \ref{muc:c4-dieuphoi}, \ref{muc:c4-congdon}, \ref{muc:c4-hople}
& Số khách bị điều phối dời giờ trả, độ dời lớn nhất, yêu cầu bị từ chối có bản ghi chặn, vế cộng
dồn, các bộ đếm kiểm tra\\
CH2. Giữ trần đổi số lượt phục vụ, số lượt đúng giờ và độ dài lời hứa thế nào so với điều phối không
trần, luật hạn chót và cách đệm lời hứa?
& \ref{muc:c4-ketqua}, \ref{muc:c4-chontran}, \ref{muc:c4-vaitro}, \ref{muc:c4-dem}
& Lượt phục vụ, lượt trễ và không trễ so với lời hứa đầu, quyết định phải phục hồi, các mức trần, thứ
tự tiêu chí của bộ chọn, độ dài thêm của lời hứa\\
CH3. Lời hứa theo hồ sơ tắc đường lịch sử đổi số lượt đúng giờ thế nào, và phần nào còn chưa được
giải thích?
& \ref{muc:c4-dubao}, \ref{muc:c4-tuyet}
& Lượt trễ khi có và không có dự báo, giá của dự báo, ba phần của sổ trên ngày thường và chiều tuyết\\
\bottomrule
\end{tabular}
\end{table}
% EV: — (lộ trình, không có khẳng định; câu hỏi theo OUTLINE-V2 cuối §1.2)

% -----------------------------------------------------------------------------------------------------
\section{Thiết lập thực nghiệm}
\label{muc:c4-thietlap}

\subsection{Kịch bản thử}
\label{muc:c4-kichban}

Mọi lần chạy dùng bộ mô phỏng FleetPy 1.0.2 với bộ dữ liệu chuẩn Manhattan của nó, nối với bộ chạy
RideBound qua bộ chuyển đổi ở mục~\ref{muc:kientruc}. Nhu cầu lấy từ các chuyến taxi của New York trong
tuần 11--18/11/2018 có sẵn trong bộ dữ liệu này~\cite{engelhardt2026}. Một kịch bản (scenario) là một tổ
hợp ngày, mẫu nhu cầu và khung giờ hai tiếng: 108 yêu cầu được lấy từ một mẫu 10\% số chuyến của khung
đó và được phục vụ bởi 8 xe 4 chỗ (Bảng~\ref{bang:c4-thamso}).
% EV: B01; B02; J18; B03

Tắc đường là một hệ số chung cho cả thành phố, đổi theo khung 15 phút và lấy từ dữ liệu chuyến taxi,
không phải tốc độ đo trên từng đoạn đường; dưới đây gọi là hệ số tắc đường đo. Mỗi cặp kịch bản và cấu
hình chạy một lần, cho một kết quả tất định. Ba khung giờ của cùng ngày và cùng mẫu trùng nhau 3 đến 20
trong 108 yêu cầu, nên các kịch bản không phải các mẫu độc lập.
% EV: B06; K49; B11; I05

\begin{table}[htbp]
\centering
\caption{Tham số chung của một kịch bản thử. Các giá trị là cấu hình của thực nghiệm, không phải chuẩn
ngành.}
\label{bang:c4-thamso}
\small
\begin{tabular}{@{}p{4.0cm}p{11.0cm}@{}}
\toprule
\textbf{Tham số} & \textbf{Giá trị}\\
\midrule
Khung giờ & 7:00--9:00, 8:00--10:00 hoặc 17:00--19:00, gọi là khung 7, 8 và 17 giờ; riêng chiều tuyết
là 14:00--16:00 và 15:00--17:00\\
Yêu cầu & 108 mỗi kịch bản, lấy từ một mẫu 10\% số chuyến taxi của khung giờ\\
Xe & 8 xe, mỗi xe 4 chỗ\\
Đồ thị & 96 nút, là bao đóng các cặp điểm đi, điểm đến lấy từ dữ liệu Manhattan, không phải mạng đường
phố\\
Ràng buộc dịch vụ & Đón trong vòng 600 giây kể từ lúc gọi; ngồi xe không quá 1,5 lần thời gian đi
thẳng (Bảng~\ref{bang:rangbuoc})\\
Thời gian dừng đón, trả & 0 giây (kiểm trên 3 lần chạy mẫu)\\
Tắc đường & Một hệ số chung cho cả thành phố theo khung 15 phút; mức nền lấy lúc 00:00 cho mọi kịch
bản\\
Bộ sinh phương án & Tối đa 100 phương án mỗi xe; mỗi lượt quyết định xét tối đa 2 yêu cầu mới\\
Bộ giải & OR-Tools 9.15, hạt giống ngẫu nhiên cố định 7\\
\bottomrule
\end{tabular}
\end{table}
% EV (khung giờ, yêu cầu): B08; B20; B02
% EV (xe, đồ thị, ràng buộc dịch vụ): B03; K48
% EV (dừng, tắc đường): B05 (chỉ 3 lần chạy mẫu, SERV:2-4; R23 #24); B06; B07
% EV (bộ sinh, bộ giải): A24; R23 #105 (giới hạn 2 yêu cầu áp cho cả lô của lượt, GEN:73-94)

\subsection{Các tập kịch bản và cấu hình được so}
\label{muc:c4-taptu}

Chương này dùng bốn tập kịch bản (Bảng~\ref{bang:c4-taptu}). Tập kiểm tra chính cho các kết quả chính:
34 trong 49 kịch bản dựng được từ các ngày 14--18/11/2018; 15 kịch bản còn lại, chia theo cặp ngày và
mẫu, dành cho hiệu chỉnh và không dùng ở đây. Ngày 14--16/11 là ngày thường, 17 và 18/11 là cuối tuần.
% EV: B08; B09; B21 (11/11/2018 là Chủ nhật, nên 14/11 là thứ Tư); B15; B17; B20

\begin{table}[htbp]
\centering
\caption{Bốn tập kịch bản dùng trong chương. Ký hiệu cấu hình theo Bảng~\ref{bang:luat}; FREE là ghép
tự do. Cột cuối là số lần chạy, mỗi lần một cặp kịch bản và cấu hình.}
\label{bang:c4-taptu}
\footnotesize
\setlength{\tabcolsep}{3pt}
\begin{tabular}{@{}p{2.5cm}p{1.9cm}p{3.1cm}p{3.0cm}p{2.6cm}p{1.9cm}@{}}
\toprule
\textbf{Tập} & \textbf{Quy mô} & \textbf{Ngày, khung giờ} & \textbf{Cấu hình dùng ở chương này} &
\textbf{Bộ chạy} & \textbf{Lần chạy}\\
\midrule
Tập kiểm tra chính & 34 kịch bản, 3.672 yêu cầu & 14--18/11/2018; khung 7, 8, 17 giờ có 9, 14, 11
kịch bản & U, C-30, C-120, M$^{+}$-30, M$^{+}$-60, V-30 & Bản lưu trữ, luật phục hồi không tệ hơn &
170 (năm cấu hình) và 34 (C-120)\\
% EV: B08; B09; B10 (170 = 34 × 5); L1a-01 (34 lần chạy C-120); A19; A28; X3-10 (run_t38c120.py:3-8)
Tập chỉnh trần buổi sáng & 8 kịch bản, 864 yêu cầu & 12--13/11/2018; khung 7 giờ & U; C và M$^{+}$ ở
30--300 giây & Như tập chính & 88 và 104\\
% EV: B15; B16; D01; D04; A28 (RUNW07:40-41); X3-10 (run_w07.py:8; run_w07x.py:7)
Tập tắc đường cố định & 16 kịch bản, 1.728 yêu cầu & 12--13/11/2018; khung 8 và 17 giờ; hệ số không
đổi trong lần chạy & FREE, U, C-30, C-60, C-120 & Bản cũ, không có luật phục hồi & 16 mỗi cấu hình\\
% EV: B17; D08; D09; H13 (n = 16 mỗi cấu hình); LAD:2-6 (travelSnapshotVersion {1: 16}: một bảng giờ đi)
Chiều tuyết 15/11/2018 & 8 kịch bản, 601 lượt phục vụ & 15/11/2018; khung 14 và 15 giờ & C-30 &
Chưa kiểm lại & 8\\
% EV: B20; N-26 (8 lần chạy C-30); notes ch.3 §5 (bộ chạy của chiều tuyết chưa kiểm)
\bottomrule
\end{tabular}
\end{table}

Điều phối không trần, luật hạn chót và luật trần độ dời khách thấy dùng cùng bộ sinh, bộ kiểm và bộ chọn
với luật C (Bảng~\ref{bang:luat}); điều phối không trần chỉ khác luật C trần 30 giây ở trần giờ trả. Ghép
tự do dùng chính sách chi phí cuốn chiếu mặc định, không qua bộ kiểm cứng. Bộ chạy và luật phục hồi của
từng tập ghi ở cột thứ năm của Bảng~\ref{bang:c4-taptu}.
% EV: A08; H01; A17; K50 (HO:109); R23 #51; B17; A19; A28; X3-10 (R23 #76, #115: OK)

Mỗi lần chạy là một cặp kịch bản và cấu hình; lần chạy hỏng được thay bằng bản chạy lại hoàn tất
(Phụ lục~\ref{phuluc:C}). Kho dữ liệu của đề tài có 1.991 bản ghi lần chạy; chương này chỉ dùng các tập ở
Bảng~\ref{bang:c4-taptu}.
% EV: ADDENDUM-1 mục 0.3 (51 thư mục C-30 = 34 + 15 + 2 bản lỗi gốc); ADD1:6 (dùng bản chạy lại); L1a-01 (34/34 đạt); B19

\subsection{Chỉ số}
\label{muc:c4-chiso}

Chỉ số đúng giờ tính so với lời hứa đầu và chỉ trên khách được phục vụ (Bảng~\ref{bang:c4-chiso}).
Ngưỡng chính là 60 giây; ngưỡng 120 giây được báo kèm vì có phép so mà hai ngưỡng cho kết luận ngược
nhau. Tỉ lệ \emph{gộp theo khách} chia tổng số khách thỏa điều kiện cho tổng số khách được phục vụ;
\emph{trung bình theo kịch bản} là trung bình không trọng số của tỉ lệ từng kịch bản. Chênh lệch là cấu
hình được thử trừ đối chứng; số dương nghĩa là nhiều hơn.
% EV: C06 (trễ so với lời hứa đầu); C08 (ở 120 s, U hơn C-30 5 lượt); C10 (đệm hòa ở 60 s, hơn ở 120 s); C27 (hai cách gộp cho hai số khác nhau); K24

Kết quả ``có dự báo'' là phép phát lại trên các lần chạy đã ghi (mục~\ref{muc:dubao}): quyết định điều
phối giữ nguyên, chỉ lời hứa đầu được thay bằng lời hứa theo dự báo. Trên tập kiểm tra chính, dự báo chỉ
có ở 21 kịch bản ngày thường; 13 kịch bản cuối tuần giữ lời hứa không dự báo. Ngày 12 và 13/11 của tập
chỉnh trần buổi sáng không có dự báo.
% EV: C09 (lưu ý: phát lại; dự báo hoạt động ở 21/34); F02; F03 (chỉ vế phát lại); N-28; D02 (lưu ý: 12–13/11 không có dự báo)

\begin{table}[htbp]
\centering
\caption{Các chỉ số dùng trong chương. $a$ là giờ trả thực tế, $p_0$ là lời hứa đầu không dự báo, $D_0$
là lời hứa đầu theo dự báo (không có dự báo thì $D_0 = p_0$); so sánh ``quá'' là so sánh chặt.}
\label{bang:c4-chiso}
\footnotesize
\setlength{\tabcolsep}{3pt}
\begin{tabular}{@{}p{2.7cm}p{5.3cm}p{3.5cm}p{3.5cm}@{}}
\toprule
\textbf{Chỉ số} & \textbf{Định nghĩa} & \textbf{Mẫu số, cách gộp} & \textbf{Vì sao đo}\\
\midrule
Lượt phục vụ & Yêu cầu được nhận và chở tới nơi & Đếm trên số yêu cầu & Giá của luật tính bằng
chuyến\\
Trễ quá $T$ & Khách được phục vụ có $a - D_0 > T$; $T$ = 60, 120, 300 hoặc 600 giây & Đếm trên khách
được phục vụ & Lời hứa bị vượt ở phía muộn\\
Không trễ quá $T$ & Lượt phục vụ trừ lượt trễ quá $T$ & Chia cho số yêu cầu; yêu cầu bị từ chối không
tính là đúng giờ & Phục vụ và đúng giờ trên một mẫu số\\
Tới sớm quá 60 giây & Khách được phục vụ có $a - D_0 < -60$ giây & Đếm trên khách được phục vụ & Lời
hứa quá dài cũng sai\\
Lời hứa dài thêm & $D_0 - p_0$ & Trung bình gộp theo khách được phục vụ & Giá của dự báo hay đệm\\
Khách đi chung & Có khách khác trên cùng xe trong khoảng từ lúc được đón tới lúc được trả & Gộp theo
khách được phục vụ & Mức ghép chuyến còn lại\\
Khách bị điều phối dời giờ trả & Có ít nhất một lần dời khác 0 do điều phối; độ dời là $\sum|c|$ &
Gộp theo khách được phục vụ & Đại lượng luật C giới hạn\\
Quyết định phải phục hồi & Lượt quyết định mà luật lời hứa loại hết phương án của một xe, kể cả phương
án giữ nguyên (mục~\ref{muc:doichung}) & Trung bình theo kịch bản; số gộp ghi riêng & Luật có tự giữ
được nghĩa vụ không\\
Khách bị chính luật phá lời hứa & Luật C: $\sum|c| > B$; luật hạn chót: $a - p_0 > B$; luật trần độ dời
khách thấy: $\sum|z| > B$ & Trung bình theo kịch bản; số gộp ghi riêng & Luật có giữ được điều nó hứa
không\\
\bottomrule
\end{tabular}
\end{table}
% EV (phục vụ): C01; L1a-01 (2.387 hoàn tất + 1.285 bị từ chối = 3.672); N-07 (khách được nhận = khách được phục vụ)
% EV (trễ, không trễ, sớm): C06; C08 (DIALPY:8-9: bị từ chối không tính là kịp); N-04 (trễ > 60 s và sớm < −60 s rời nhau); C09; C10
% EV (dài thêm, đi chung, bị dời): F06; N-06 (gộp theo khách); N-09 (định nghĩa đi chung); C20; N-07; A06
% EV (phải phục hồi, phá lời hứa): C13 (trung bình theo ô); N-21 (định nghĩa phá hứa của M⁺-30 và C-30; số gộp); A11; GLOSSARY-VI mục E

\subsection{Mức trần 30 giây và tính chất thăm dò của kết quả}
\label{muc:c4-thamdo}

Mức trần chính $B$ = 30 giây là mức đã dùng cho cấu hình có trần ở tập kịch bản tĩnh trước đó
(Phụ lục~\ref{phuluc:A}), được giữ nguyên cho tập kiểm tra chính và không chỉnh theo kết quả của tập này.
Mức 30 giây không phải ngưỡng cảm nhận của khách hay chuẩn ngành, nên mỗi mức trần được trình bày như một
điểm trên đường đánh đổi (mục~\ref{muc:c4-chontran}).
% EV: K33; X3-18 (TEX/08-chuong4.tex:128-136: C1 = trần 30 giây); B10 (năm cấu hình cố định trước khi chạy tập chính; không nêu kế hoạch); RD LR-31; D20; K52; K36

Các kết quả của chương là kết quả thăm dò. Ngưỡng phân tích, cách chia tập và mức thử 120 giây được chọn
sau khi đã xem dữ liệu của các tập phát triển, và một số kịch bản đã được xem khi chạy thử hệ thống. Các
con số mô tả những gì đã đo trên các kịch bản này, không phải phép kiểm trên dữ liệu chưa từng xem.
% EV: I06; B08 (lưu ý: d20181114-s10-r1 đã xem ở lần chạy thử, T38:95); B10 (lưu ý: ngưỡng, biên, cách chia chọn sau khi xem); K42

% -----------------------------------------------------------------------------------------------------
\section{Kết quả chính khi chưa có dự báo}
\label{muc:c4-ketqua}

Trên 3.672 yêu cầu của tập kiểm tra chính, luật C với trần 30 giây phục vụ 2.338 lượt, ít hơn điều phối
không trần 96 lượt. Số lượt trễ quá một phút so với lời hứa đầu giảm từ 203 xuống 76, nên số lượt không
trễ quá một phút tăng từ 2.231 lên 2.262.
% EV: B09; C01; C02; C06; C08

Luật hạn chót ở đây chỉ mượn phần hạn chót của Schulz và Pfeiffer~\cite{schulz2026}, ở mức 30 và 60 giây
thay vì 900 giây mặc định của họ. Nó chạy cùng bộ chạy, bộ sinh, bộ chọn và luật phục hồi không tệ hơn
như luật C. Vì vậy đây là phép so hai luật lời hứa trong cùng một bộ điều phối, không phải phép so với hệ
thống của hai tác giả.
% EV: J10; J11 (tập chính 30 và 60 s; R23 #57); A09; A19; H01; RD CX-07, LR-05

\begin{table}[htbp]
\centering
\caption{Kết quả trên tập kiểm tra chính với lời hứa đầu không dự báo (34 kịch bản, 3.672 yêu cầu mỗi
cấu hình). Số đếm gộp trên mọi kịch bản; trễ tính so với lời hứa đầu; hai cột ``không trễ quá'' tính trên
3.672 yêu cầu.}
\label{bang:c4-chinh}
\small
\begin{tabular}{@{}lrrrrr@{}}
\toprule
\textbf{Cấu hình} & \textbf{Phục vụ} & \textbf{Trễ quá 60 s} & \textbf{Không trễ quá 60 s} &
\textbf{Không trễ quá 120 s} & \textbf{Khách đi chung}\\
\midrule
U & 2.434 & 203 & 2.231 & 2.332 & 555\\
C-30 & 2.338 & 76 & 2.262 & 2.327 & 313\\
M$^{+}$-30 & 2.338 & 74 & 2.264 & 2.327 & 307\\
% EV (U, C-30, M⁺-30): C01; C06; C08; N-09 (555 / 313 / 307, số đếm chính xác)
M$^{+}$-60 & 2.347 & \multicolumn{4}{c}{chưa đưa vào bảng}\\
V-30 & 2.336 & \multicolumn{4}{c}{chưa đưa vào bảng}\\
% EV (M⁺-60, V-30): C01 (chỉ số phục vụ có trong sổ)
\bottomrule
\end{tabular}
\end{table}

Cách đọc bảng: dưới trần 30 giây, 2.338 yêu cầu được phục vụ; 76 lượt trong số đó tới muộn hơn lời hứa
đầu quá 60 giây, nên 2.338 $-$ 76 = 2.262 lượt không trễ quá 60 giây.
% EV: C01; C06; C08

Cái giá của trần trước hết là số lượt phục vụ. Mức giảm 96 lượt bằng 2,61 điểm phần trăm số yêu cầu,
hay 3,94\% số lượt của điều phối không trần. Nó nằm chủ yếu ở khung 8 giờ (từ 1.014 xuống 975 lượt) và
17 giờ (từ 904 xuống 852); khung 7 giờ chỉ giảm 5 lượt. Số khách đi chung giảm mạnh hơn, từ 555 xuống
313, tức từ 22,80\% xuống 13,39\% số khách được phục vụ (gộp theo khách). Khóa luận không có bằng chứng
về giá trị của các chuyến đi chung này đối với khách hay hãng xe.
% EV: C02; C03; N-09; N-20 (lưu ý: D20, I07, I10)

Đổi lại, số lượt trễ quá một phút giảm 127, nhiều hơn mức giảm phục vụ, nên số lượt không trễ quá một
phút tăng 31: $+31 = -96 - (-127)$. Ở ngưỡng hai phút thì ngược lại: điều phối không trần có 2.332 lượt
không trễ quá 120 giây, hơn luật C 5 lượt. Hai cấu hình cũng không phục vụ cùng những khách: 2.207 khách
chung, 227 khách chỉ được điều phối không trần phục vụ và 131 khách chỉ được luật C phục vụ. Vì vậy 31
lượt kể trên là chênh lệch giữa hai tổng, không phải 31 khách cụ thể được đưa tới đúng giờ. Với 2.207
khách chung, luật C chờ lâu hơn 0,5 giây và ngồi xe ngắn hơn 10,5 giây, nhưng đây là tập khách hình thành
sau khi áp luật, không phải một phép so ngẫu nhiên.
% EV: C06; C08 (timely = phục vụ − trễ; ở 120 s U hơn C-30 5 lượt); C04; C05 (lưu ý: tập sau điều trị); RD CX-08

So với luật hạn chót 30 giây, luật C phục vụ đúng bằng số lượt (2.338) nhưng không phải cùng mọi khách:
2.332 khách chung, mỗi bên 6 khách riêng; số lượt trễ quá một phút là 76 so với 74. Luật hạn chót 60 giây
phục vụ 2.347 lượt, luật trần độ dời khách thấy 2.336 lượt. Như vậy, số lượt phục vụ của các luật lời hứa
gần như ngang nhau; chỗ khác nhau là mỗi luật có giữ được nghĩa vụ của chính nó hay không
(Bảng~\ref{bang:c4-nghiavu}).
% EV: C01; C04; C06; K05; K07; C32

\begin{table}[htbp]
\centering
\caption{Nghĩa vụ của chính từng luật trên tập kiểm tra chính. Hai cột tỉ lệ là trung bình theo kịch
bản (34 kịch bản); cột cuối đếm số kịch bản. Định nghĩa ``phá lời hứa'' của từng luật ở
Bảng~\ref{bang:c4-chiso}.}
\label{bang:c4-nghiavu}
\small
\begin{tabular}{@{}lrrr@{}}
\toprule
\textbf{Cấu hình} & \textbf{\begin{tabular}[b]{@{}r@{}}Quyết định phải\\ phục hồi (\%)\end{tabular}} &
\textbf{\begin{tabular}[b]{@{}r@{}}Khách bị chính luật\\ phá lời hứa (\%)\end{tabular}} &
\textbf{\begin{tabular}[b]{@{}r@{}}Kịch bản có ít nhất\\ một khách bị phá\end{tabular}}\\
\midrule
C-30 & 0 & 0 & 0\\
M$^{+}$-30 & 41,3 & 14,5 & 28\\
M$^{+}$-60 & 22,5 & 4,2 & 22\\
V-30 & 53,8 & 21,6 & 34\\
\bottomrule
\end{tabular}
\end{table}
% EV: C13 (trung bình theo ô; ô có ≥ 1 phá hứa 0/28/22/34); N-21 (C-30 0/34 kịch bản, 0 quyết định phải phục hồi)

Cách đọc bảng: với luật hạn chót 30 giây, trung bình mỗi kịch bản có 41,3\% lượt quyết định phải phục
hồi và 14,5\% khách được phục vụ tới muộn hơn lời hứa đầu quá 30 giây. Có 28 trong 34 kịch bản có ít
nhất một khách như vậy.
% EV: C13; N-21

Luật C không có quyết định nào phải phục hồi và không vượt trần của khách nào, ở cả 34 kịch bản. Điều
này đúng theo cấu tạo: phương án giữ nguyên kế hoạch có phần dời do điều phối bằng 0 nên luôn qua luật C
(Nhận xét~\ref{nx:giunguyen}). Luật hạn chót so giờ trả mới với lời hứa đầu, nên riêng phần thay đổi khi
giữ nguyên kế hoạch cũng có thể làm nó loại cả phương án giữ nguyên. Tính gộp, luật hạn chót 30 giây phải
phục hồi ở 7.710 trong 20.103 lượt quyết định (38,35\%) và phá lời hứa với 289 trong 2.338 khách
(12,4\%). Gánh nặng dồn vào khung 7 giờ: cả 9 kịch bản có khách bị phá lời hứa, trung bình 36,5\% khách
bị phá và 75,6\% lượt quyết định phải phục hồi; ở khung 17 giờ chỉ 5 trong 11 kịch bản, với 1,0\% và
5,0\%. Ở mức 60 giây, các tỉ lệ này giảm nhưng không về 0.
% EV: A21; A22; N-21 (gộp 7.710/20.103 = 38,35%; 289/2.338 = 12,36%; khung 7 giờ 9/9, 17 giờ 5/11); A10; C13; C14

Bảng~\ref{bang:c4-nghiavu} không so hai chuẩn đúng giờ: luật C hứa về phần dời do điều phối, còn luật
hạn chót hứa về toàn bộ độ trễ so với lời hứa đầu. Điều bảng cho thấy là luật nào giữ được điều nó hứa
trên các kịch bản đã chạy.
% EV: C13 (lưu ý: hai nghĩa vụ khác nhau); K05

Bảng~\ref{bang:c4-phanbotre} cho thấy số lượt trễ giảm ở mọi ngưỡng, không chỉ ở ngưỡng một phút.
% EV: N-24

\begin{table}[htbp]
\centering
\caption{Phân bố trễ so với lời hứa đầu không dự báo trên tập kiểm tra chính: số lượt trễ quá mỗi ngưỡng
(trong số khách được phục vụ) và số lượt đó tính trên 100 yêu cầu (3.672 yêu cầu mỗi cấu hình).}
\label{bang:c4-phanbotre}
\small
\begin{tabular}{@{}lrrrr@{}}
\toprule
& \multicolumn{2}{c}{\textbf{U}} & \multicolumn{2}{c}{\textbf{C-30}}\\
\cmidrule(lr){2-3}\cmidrule(lr){4-5}
\textbf{Trễ quá} & \textbf{Lượt} & \textbf{Trên 100 yêu cầu} & \textbf{Lượt} & \textbf{Trên 100 yêu cầu}\\
\midrule
1 phút & 203 & 5,53 & 76 & 2,07\\
2 phút & 102 & 2,78 & 11 & 0,30\\
5 phút & 18 & 0,49 & 0 & 0\\
10 phút & 2 & 0,05 & 0 & 0\\
\midrule
Trễ lớn nhất & \multicolumn{2}{c}{11 phút} & \multicolumn{2}{c}{3,87 phút}\\
\bottomrule
\end{tabular}
\end{table}
% EV: N-24 (203/102/18/2 và 76/11/0/0; 5,53/2,78/0,49/0,05 và 2,07/0,30/0/0); C07; C28; C26 (lâu nhất 11 phút và 3,87 phút)

Cách đọc bảng: khi không trần, 18 lượt tới muộn hơn lời hứa đầu quá 5 phút, tức 0,49 lượt trên 100 yêu
cầu.
% EV: N-24; C28

Khi không trần, 2 lượt trễ quá 10 phút và lượt trễ nhất muộn khoảng 11 phút; dưới trần 30 giây, không
lượt nào trễ quá 5 phút và lượt trễ nhất muộn 3,87 phút. Trong 76 khách còn trễ quá một phút dưới luật C,
72 người không có lần dời giờ trả do điều phối nào được ghi. Phần trễ này là phần mà lời hứa theo dự báo
nhắm tới (mục~\ref{muc:c4-dubao}).
% EV: N-24; C26; C21 (4/76 từng bị dời); K39

Để xem kết quả có dựa vào vài kịch bản hay không, Bảng~\ref{bang:c4-ngaykhung} tách chênh lệch theo ngày
và khung giờ.

\begin{table}[htbp]
\centering
\caption{Chênh lệch luật C trần 30 giây trừ điều phối không trần theo ngày và khung giờ, lời hứa đầu
không dự báo. Số đếm cộng trên các kịch bản của mỗi nhóm. Cột cuối là số kịch bản mà luật C có số lượt
không trễ quá 60 giây không ít hơn.}
\label{bang:c4-ngaykhung}
\small
\begin{tabular}{@{}lrrrrr@{}}
\toprule
\textbf{Ngày, khung} & \textbf{Kịch bản} & \textbf{Phục vụ} & \textbf{Trễ quá 60 s} &
\textbf{Không trễ quá 60 s} & \textbf{Kịch bản C không kém}\\
\midrule
14/11, 7 giờ & 2 & 0 & 0 & 0 & 2\\
14/11, 8 giờ & 2 & $-4$ & $-7$ & $+3$ & 2\\
14/11, 17 giờ & 2 & $-8$ & $-7$ & $-1$ & 1\\
15/11, 7 giờ & 3 & $+2$ & $-2$ & $+4$ & 3\\
15/11, 8 giờ & 3 & $-8$ & $-10$ & $+2$ & 3\\
16/11, 7 giờ & 3 & $-3$ & $-1$ & $-2$ & 2\\
16/11, 8 giờ & 3 & $-9$ & $-13$ & $+4$ & 3\\
16/11, 17 giờ & 3 & $-11$ & $-21$ & $+10$ & 2\\
17/11, 7 giờ & 1 & $-4$ & $-6$ & $+2$ & 1\\
17/11, 8 giờ & 3 & $-2$ & $-7$ & $+5$ & 3\\
17/11, 17 giờ & 3 & $-19$ & $-18$ & $-1$ & 2\\
18/11, 8 giờ & 3 & $-16$ & $-15$ & $-1$ & 2\\
18/11, 17 giờ & 3 & $-14$ & $-20$ & $+6$ & 3\\
\midrule
Cộng & 34 & $-96$ & $-127$ & $+31$ & 29\\
\bottomrule
\end{tabular}
\end{table}
% EV: N-22 (bảng dữ liệu cho Bảng 4.7, ADD1:167-186; cột phục vụ tổng −96 = 2.338 − 2.434, ADD1:12,17)

Cách đọc bảng: ngày 16/11, khung 17 giờ, trên 3 kịch bản, luật C phục vụ ít hơn 11 lượt và có ít hơn 21
lượt trễ quá 60 giây, nên có thêm 10 lượt không trễ quá 60 giây; luật C không kém ở 2 trong 3 kịch bản.
% EV: N-22

Ở cả 34 kịch bản, luật C không có nhiều lượt trễ quá một phút hơn điều phối không trần, và ở 29 kịch bản
số lượt không trễ quá một phút của luật C không ít hơn. Ở khung 8 và 17 giờ, điều phối không trần có
nhiều lượt trễ hơn ở cả 25 kịch bản; ở khung 7 giờ, nhiều hơn ở 3 kịch bản và bằng nhau ở 6. Vì 34 kịch
bản thuộc 5 ngày và không độc lập với nhau, các phép đếm theo kịch bản này là mô tả, không phải kiểm định
thống kê.
% EV: N-22 (34/34 trễ ≤ U; 29/34 không trễ ≥ U); C30; I05; B11

% -----------------------------------------------------------------------------------------------------
\section{Phần do điều phối: luật C chặn gì và mất gì}
\label{muc:c4-dieuphoi}

Dưới luật C với trần 30 giây, 33 trong 2.338 khách được phục vụ có giờ trả bị điều phối dời, mỗi khách
nhiều nhất 29,2 giây. Khi không trần, con số này là 174 trong 2.434 khách, và một khách bị dời tổng cộng
590,1 giây (gần 10 phút).
% EV: C20; N-23

Bảng~\ref{bang:c4-doidieuphoi} so phần dời do điều phối giữa điều phối không trần và hai mức trần trên
tập kiểm tra chính; cột trần 120 giây được dùng lại ở mục~\ref{muc:c4-chontran}.

\begin{table}[htbp]
\centering
\caption{Phần dời giờ trả do điều phối trên tập kiểm tra chính, lời hứa đầu không dự báo. Mẫu số của
dòng đầu là số khách được phục vụ (2.434, 2.338 và 2.387); độ dời của một khách là tổng $|c|$ của khách
đó.}
\label{bang:c4-doidieuphoi}
\small
\begin{tabular}{@{}p{7.6cm}rrr@{}}
\toprule
\textbf{Chỉ số} & \textbf{U} & \textbf{C-30} & \textbf{C-120}\\
\midrule
Khách bị điều phối dời giờ trả & 174 (7,15\%) & 33 (1,41\%) & 103 (4,32\%)\\
Khách bị dời từ hai lần trở lên & 12 & 0 & 0\\
Tổng dời lớn nhất của một khách (s) & 590,1 & 29,2 & 118,9\\
Trung vị độ lớn mỗi lần dời (s) & 108,6 & 8,3 & 62,5\\
% EV: C20; C22 (trung vị theo lần: 108,6 s trên 187 lần, 8,3 s trên 33 lần); N-07 (C-120: 103, 0 khách ≥ 2 lần, nên trung vị theo lần = trung vị theo khách 62,54 s); N-23
Khách trễ quá 60 s từng bị dời & 134 trên 203 & 4 trên 76 & 60 trên 132\\
Theo sổ: phần tính cho điều phối trong số giây trễ của khách trễ quá 60 s (\%) & 73,4 & 1,2 & 38,8\\
% EV: C21; C24; N-23 (cột C-120: 60/132, 38,8%)
\bottomrule
\end{tabular}
\end{table}

Cách đọc bảng: khi không trần, 174 khách bị điều phối dời giờ trả ít nhất một lần, 12 người trong số đó
bị dời từ hai lần, và tổng dời lớn nhất của một khách là 590,1 giây.
% EV: C20

Dưới trần 30 giây, mỗi khách bị dời chỉ bị dời một lần, trung vị 8,3 giây mỗi lần; khi không trần, trung
vị là 108,6 giây trên 187 lần dời. Trong 33 khách bị dời dưới trần 30 giây, 29 người bị dời muộn đi; 4
người còn lại có giờ trả bị kéo sớm hơn, mỗi người đúng 1 mili giây. Khi không trần, 134 trong 203 khách
trễ quá một phút từng bị điều phối dời giờ trả; con số này là giao của hai tập khách, không phải tỉ lệ
nhân quả. Theo sổ theo từng khách (công thức~\eqref{eq:so}), phần tính cho điều phối bằng 73,4\% tổng số
giây trễ của 203 khách đó; dưới trần 30 giây, phần này là 1,2\%. Nói cách khác, dưới trần 30 giây,
gần như toàn bộ số giây trễ còn lại nằm ở hai phần kia của sổ.
% EV: E02; C22; C23 (29 muộn, 4 sớm); N-05 (mỗi lần dời sớm đúng −1 ms; chỉ kiểm ở U, C-30, C-120, M⁺-30 trên 34 kịch bản); C21 (lưu ý: giao hai tập); C24 (lưu ý: sổ mô hình, không nhân quả); K04

Trần vẫn giữ ở những kịch bản không có dự báo và ở ngày có sự kiện. Ở 13 kịch bản cuối tuần, nơi không có
dự báo, tổng dời lớn nhất của một khách là 29,2 giây với trần 30 giây và 116,5 giây với trần 120 giây. Trên 8
kịch bản chiều tuyết, dưới trần 30 giây, 8 trong 601 khách bị điều phối dời giờ trả, lớn nhất 25,2 giây,
không khách nào vượt trần và không có quyết định nào phải phục hồi trong 5.067 lượt quyết định.
% EV: N-28; N-26 (thay anchor F13 của RD §8 mục 4, ADDENDUM-1 mục 0.4)

Một ví dụ trên tập chỉnh trần buổi sáng cho thấy luật C khác luật hạn chót ở đâu. Một khách đang trên
xe đã muộn hơn lời hứa đầu 79,0 giây, hoàn toàn do phần thay đổi khi giữ nguyên kế hoạch; phần dời do
điều phối của khách lúc đó bằng 0. Một phương án ghép thêm một khách mới sẽ dời giờ trả của khách này
thêm 68,3 giây. Với mức 90 giây, luật hạn chót loại phương án này vì giờ trả mới muộn hơn lời hứa đầu
khoảng 147 giây; luật C cho qua vì phần dời do điều phối chỉ là 68,3 giây. Ba chỗ hai luật tách nhau trên
tập này, ở mức 90 và 120 giây, đều có cùng dạng: phần thay đổi khi giữ nguyên kế hoạch 61--79 giây, phần do
điều phối 61--68 giây. Mỗi lần như vậy luật C nhận thêm đúng một khách. Đây là phép
tính theo mô hình tại cùng một trạng thái, không phải phép đo nhân quả.
% EV: D05 (79,0 + 68,3 s; 60,9 + 61,2 s; W07XM:2-16; W07X §4: 147 s > 90 s); D05 lưu ý (phép chiếu mô hình); K04

Cái giá của trần là các yêu cầu không được nhận. Dưới trần 30 giây, 1.334 yêu cầu bị từ chối. Trong số
đó, 148 yêu cầu cuối cùng bị từ chối và có ít nhất một phương án mang bản ghi chặn do trần giờ trả của
một khách đã nhận; điều phối không trần phục vụ 116 trong 148 yêu cầu này. 1.186 yêu cầu bị từ chối còn
lại chỉ có lý do cấu trúc, không có bản ghi chặn do trần. Bản ghi chặn cũng không phải là khách: 845 bản
ghi của trần 30 giây trên tập chính nằm ở 643 yêu cầu, và 495 yêu cầu trong số đó cuối cùng vẫn được
phục vụ bằng phương án khác.
% EV: D13 (148; 116 U chở; 1.186; tổng 1.334); K10; D16 (643 cặp, 495 phục vụ, 148 từ chối); E01 (845)

Với 148 yêu cầu bị từ chối đó, mức vượt trần nhỏ nhất được ghi có trung vị 109,6 giây và lớn nhất 524,8
giây. Chỉ 2 trong 148 khách chặn đã dùng phần trần trước đó, và ở cả 148 yêu cầu, bản ghi đều có một lần
dời đơn lẻ vượt 30 giây. Ở mức trần 120 giây, ngưỡng được ghi của 59 trong 148 yêu cầu qua được phép sàng
lọc. Vì bộ kiểm dừng ở lỗi gặp trước và quỹ đạo sẽ đổi khi trần đổi, điều đó không có nghĩa 59 yêu cầu sẽ
được phục vụ; mục~\ref{muc:c4-chontran} cho kết quả chạy thật ở mức này. Các phân loại trong mục này
dựa vào bản ghi dừng ở lỗi gặp trước, nên không chứng minh một yêu cầu bị từ chối chỉ vì trần.
% EV: D14 (trung vị 109,6 s; lớn nhất 524,8 s; 11/59/135 ở 60/120/300 s); D19 (2/148); K11; A23; I08; D13 (lưu ý)

% -----------------------------------------------------------------------------------------------------
\section{Chọn mức trần}
\label{muc:c4-chontran}

Nới trần không cho một đường đơn điệu. Trên 3.672 yêu cầu của tập kiểm tra chính, trần 120 giây phục vụ
2.387 lượt, nằm giữa trần 30 giây (2.338) và không trần (2.434), nhưng có 2.366 lượt không trễ quá hai
phút, nhiều hơn cả hai (2.327 và 2.332).
% EV: N-06; L1a-02; C08; K12

Mức trần được đo trên ba tập, mỗi tập một câu hỏi, và kết quả của các tập không được cộng gộp;
Hình~\ref{hinh:c4-tran} đặt chúng cạnh nhau.
% EV: B17; B15; B16; D01; L1a-01; I12

\textbf{Tập tắc đường cố định.} Bảng~\ref{bang:c4-codinh} cho thang từ ghép tự do tới trần 30 giây khi
thời gian đi không đổi.

\begin{table}[htbp]
\centering
\caption{Các mức trần trên tập tắc đường cố định (16 kịch bản, 1.728 yêu cầu mỗi cấu hình; bộ chạy cũ,
không có luật phục hồi). Ba cột số đầu là số đếm; hai cột cuối là tỉ lệ gộp theo khách được phục vụ.}
\label{bang:c4-codinh}
\small
\begin{tabular}{@{}lrrrrr@{}}
\toprule
\textbf{Cấu hình} & \textbf{Phục vụ} & \textbf{Trễ quá 60 s} & \textbf{Không trễ quá 120 s} &
\textbf{Đi chung (\%)} & \textbf{Bị dời giờ trả (\%)}\\
\midrule
FREE & 1.629 & 274 & 1.465 & 53,1 & 23,8\\
U & 1.584 & 47 & 1.557 & 21,8 & 3,9\\
C-120 & 1.563 & 23 & 1.563 & 19,4 & 2,6\\
C-60 & 1.540 & 0 & 1.540 & 17,3 & 1,3\\
C-30 & 1.534 & 0 & 1.534 & 15,7 & 0,4\\
\bottomrule
\end{tabular}
\end{table}
% EV: D09 (số đếm, WP17); D08 (tỉ lệ đi chung, bị dời; gộp theo khách); B17

Cách đọc bảng: với trần 120 giây, 1.563 trong 1.728 yêu cầu được phục vụ, 23 lượt trễ quá 60 giây; 19,4\%
khách được phục vụ đi chung và 2,6\% bị điều phối dời giờ trả.
% EV: D08; D09

Siết dần từ ghép tự do tới trần 30 giây, số lượt phục vụ giảm từ 1.629 xuống 1.534 và số lượt trễ quá 60
giây từ 274 xuống 0. Điểm gãy nằm giữa 60 và 120 giây: số lượt trễ quá 60 giây tăng từ 0 lên 23. Tập này
chạy bằng bộ chạy cũ không có luật phục hồi, nên các con số không chuyển sang tập kiểm tra chính.
% EV: D09; D11; B17; I12

\textbf{Tập chỉnh trần buổi sáng.} Trên 864 yêu cầu, số lượt phục vụ chỉ dao động trong 6 lượt: 660 ở trần
30 giây, 661 ở 60 giây, 664 ở 90 giây, 666 ở 120--180 giây, 665 ở 240, 300 giây và khi không trần.
Bảng~\ref{bang:c4-sang} cho các chỉ số đúng giờ ở bốn mức.
% EV: D01; B15; B16

\begin{table}[htbp]
\centering
\caption{Bốn mức trên tập chỉnh trần buổi sáng (8 kịch bản khung 7 giờ ngày 12--13/11/2018, 864 yêu cầu
mỗi cấu hình; lời hứa đầu không dự báo). Số đếm gộp; hai cột ``không trễ quá'' tính trên 864 yêu cầu.}
\label{bang:c4-sang}
\small
\begin{tabular}{@{}lrrrr@{}}
\toprule
\textbf{Cấu hình} & \textbf{Phục vụ} & \textbf{Trễ quá 60 s} & \textbf{Không trễ quá 60 s} &
\textbf{Không trễ quá 120 s}\\
\midrule
U & 665 & 50 & 615 & 656\\
C-30 & 660 & 47 & 613 & 658\\
C-90 & 664 & 50 & 614 & 660\\
C-120 & 666 & 51 & 615 & 659\\
\bottomrule
\end{tabular}
\end{table}
% EV: D01; D02 (CAP:3-11; DIAL17:3-11, hai bộ đọc độc lập); D02 lưu ý (cột dự báo bằng ngây thơ ở 12–13/11)

Cách đọc bảng: với trần 120 giây, 666 yêu cầu được phục vụ, 51 lượt trễ quá 60 giây, 615 lượt không trễ
quá 60 giây và 659 lượt không trễ quá 120 giây.
% EV: D01; D02

Trần 120 giây phục vụ nhiều hơn điều phối không trần một lượt, nên điều phối không trần là một điểm đo,
không phải cận trên. Từ 30 lên 120 giây, trần thêm 6 lượt phục vụ, 4 lượt trễ quá 60 giây và 2 lượt
không trễ quá 60 giây. Trần lỏng cho phép phần dời lớn: ở một kịch bản của tập này, với trần 300 giây,
khách trễ nhất muộn 307,3 giây, trong đó 208,2 giây là phần do điều phối.
% EV: D01; K12; D03 (+6 / +4 / +2); N-30 (W07:49,53); D06

Luật hạn chót trả giá ở các mức chặt. Tỉ lệ lượt quyết định nó phải phục hồi (trung bình theo kịch bản)
là 58,6; 35,5; 13,0; 6,6; 6,8; 2,8 và 1,7\% ở các mức 30, 60, 90, 120, 150, 180 và 240--300 giây, còn
luật C luôn 0\%. Ở mọi mức, hai luật phục vụ chênh nhau không quá 2 lượt.
% EV: D04 (trung bình theo ô); R23 #57 (luật hạn chót chạy 30–300 s ở tập này)

\textbf{Tập kiểm tra chính.} Bảng~\ref{bang:c4-tran120} đặt trần 120 giây cạnh hai cấu hình của
mục~\ref{muc:c4-ketqua}, cả khi không có và khi có dự báo.

\begin{table}[htbp]
\centering
\caption{Trần 120 giây trên tập kiểm tra chính (34 kịch bản, 3.672 yêu cầu mỗi cấu hình). Số đếm gộp; các
dòng ``không trễ quá'' tính trên 3.672 yêu cầu. Khối dưới là phép phát lại lời hứa theo dự báo trên cùng
các lần chạy; dự báo chỉ có ở 21 kịch bản ngày thường.}
\label{bang:c4-tran120}
\small
\begin{tabular}{@{}p{8.4cm}rrr@{}}
\toprule
\textbf{Chỉ số} & \textbf{U} & \textbf{C-30} & \textbf{C-120}\\
\midrule
\multicolumn{4}{@{}l}{\textit{Lời hứa đầu không dự báo}}\\
Lượt phục vụ & 2.434 & 2.338 & 2.387\\
Lượt trễ quá 60 s & 203 & 76 & 132\\
Lượt trễ quá 120 s & 102 & 11 & 21\\
Lượt trễ quá 300 s & 18 & 0 & 0\\
Không trễ quá 60 s & 2.231 & 2.262 & 2.255\\
Không trễ quá 120 s & 2.332 & 2.327 & 2.366\\
Khách đi chung & 555 & 313 & 428\\
% EV: L1a-02; L1a-04 (U, C-30 tái lập); N-06; C06; C07; C08; N-09 (L1a-07)
\midrule
\multicolumn{4}{@{}l}{\textit{Lời hứa đầu theo dự báo}}\\
Lượt trễ quá 60 s & 137 & 10 & 65\\
Không trễ quá 60 s & 2.297 & 2.328 & 2.322\\
Không trễ quá 120 s & 2.347 & 2.338 & 2.379\\
Lời hứa dài thêm, trung bình gộp theo khách (s) & 7,18 & 7,37 & 7,27\\
% EV: L1a-03; N-03; N-06; N-16; C09; F06
\bottomrule
\end{tabular}
\end{table}

Cách đọc bảng: với trần 120 giây và lời hứa đầu không dự báo, 2.387 yêu cầu được phục vụ, 132 lượt trễ
quá 60 giây, nên 2.255 lượt không trễ quá 60 giây trên 3.672 yêu cầu.
% EV: N-06; L1a-02

So với trần 30 giây, trần 120 giây phục vụ thêm 49 lượt và có thêm 56 lượt trễ quá một phút, nên ít hơn 7
lượt không trễ quá một phút. Ngược lại, nó thêm 39 lượt không trễ quá hai phút và ở ngưỡng này cao hơn cả
hai điểm đầu mút đã đo. Khi có dự báo, trần 120 giây có 2.322 lượt không trễ quá một phút và 2.379 lượt
không trễ quá hai phút, tức ít hơn trần 30 giây 6 lượt và thêm 41 lượt. Độ dài thêm của lời hứa gần như
bằng nhau (7,27 so với 7,37 giây). Mục~\ref{muc:c4-dem} đặt điểm này cạnh cách đệm lời hứa.
% EV: L1a-09 (+49 / +56 / −7 / +39; dự báo −6 / +41); N-06; L1a-02 (2.366 > U 2.332 và C-30 2.327); K12; D01; L1a-03; N-16

Dưới trần 120 giây, 9 khách dùng từ 90\% trần trở lên, không khách nào vượt trần và không có quyết định
nào phải phục hồi trong 20.435 lượt quyết định; số khách bị dời và phần tính cho điều phối ở cột cuối của
Bảng~\ref{bang:c4-doidieuphoi}.
% EV: N-07; L1a-05; N-23

Hai tập khách được phục vụ khác nhau: trần 120 giây và trần 30 giây chung 2.266 lượt, riêng 121 và 72
lượt, và 93 trong 121 lượt đó cũng được phục vụ khi không trần. Vì vậy 49 lượt phục vụ thêm là chênh lệch
hai tổng, không phải 49 trong 96 lượt mà trần 30 giây bỏ so với không trần. Phần phục vụ thêm và phần trễ
thêm dồn vào khung 8 giờ (1.001 so với 975 lượt phục vụ; 49 so với 15 lượt trễ quá một phút) và 17 giờ
(871 so với 852; 20 so với 0). So từng kịch bản, trần 120 giây phục vụ nhiều hơn ở 23, bằng ở 8 và ít
hơn ở 3.
% EV: N-10 (lưu ý: không viết "lấy lại 49 trong 96 lượt"); L1a-08; N-11; L1a-09

\begin{figure}[htbp]
\centering
\begin{tikzpicture}[font=\footnotesize,
  vuong/.style={draw,fill=white,rectangle,minimum size=4.2pt,inner sep=0pt},
  tamgiac/.style={draw,fill=white,regular polygon,regular polygon sides=3,minimum size=6pt,inner sep=0pt},
  tron/.style={draw,fill=white,circle,minimum size=4.6pt,inner sep=0pt},
  thoi/.style={draw,fill=black,diamond,minimum size=6.5pt,inner sep=0pt},
  nhan/.style={font=\scriptsize,inner sep=1.5pt}]
% (a) Tập tắc đường cố định. Trục ngang: phục vụ − 1.530, 0,06 cm mỗi lượt; trục dọc: trễ, 0,05 cm mỗi lượt.
\begin{scope}[shift={(0,0)}]
\draw[->] (0,0)--(3.95,0);
\draw[->] (0,0)--(0,3.35);
\foreach \vx/\so in {0/1.530,1.2/1.550,2.4/1.570,3.6/1.590}
  \draw (\vx,0)--(\vx,-0.07) node[below,nhan]{\so};
\foreach \vy/\so in {0/0,1/20,2/40,3/60}
  \draw (0,\vy)--(-0.07,\vy) node[left,nhan]{\so};
\draw[dashed] (0.24,0)--(0.60,0)--(1.98,1.15)--(3.24,2.35);
\node[vuong] at (0.24,0) {}; \node[nhan,above] at (0.24,0.05) {30};
\node[vuong] at (0.60,0) {}; \node[nhan,above] at (0.62,0.05) {60};
\node[vuong] at (1.98,1.15) {}; \node[nhan,above left] at (1.98,1.15) {120};
\node[vuong] at (3.24,2.35) {}; \node[nhan,above left] at (3.24,2.35) {U};
\node at (1.8,-0.72) {Lượt phục vụ};
\node[rotate=90] at (-0.85,1.6) {Lượt trễ quá 60 s};
\node[anchor=south] at (1.9,3.4) {(a) Tắc đường cố định};
\end{scope}
% (b) Tập chỉnh trần buổi sáng. Trục ngang: phục vụ − 658, 0,36 cm mỗi lượt; trục dọc: trễ, 0,05 cm mỗi lượt.
\begin{scope}[shift={(5.3,0)}]
\draw[->] (0,0)--(3.95,0);
\draw[->] (0,0)--(0,3.35);
\foreach \vx/\so in {0/658,0.72/660,1.44/662,2.16/664,2.88/666,3.6/668}
  \draw (\vx,0)--(\vx,-0.07) node[below,nhan]{\so};
\foreach \vy/\so in {0/0,1/20,2/40,3/60}
  \draw (0,\vy)--(-0.07,\vy) node[left,nhan]{\so};
\draw[densely dotted,thick] (0.72,2.35)--(2.16,2.50)--(2.88,2.55)--(2.52,2.50);
\node[tamgiac] at (0.72,2.35) {}; \node[nhan,below] at (0.72,2.25) {30};
\node[tamgiac] at (2.16,2.50) {}; \node[nhan,below] at (2.16,2.40) {90};
\node[tamgiac] at (2.88,2.55) {}; \node[nhan,right] at (2.98,2.55) {120};
\node[tamgiac] at (2.52,2.50) {}; \node[nhan,above] at (2.52,2.60) {U};
\node at (1.8,-0.72) {Lượt phục vụ};
\node[anchor=south] at (1.9,3.4) {(b) Buổi sáng};
\end{scope}
% (c) Tập kiểm tra chính. Trục ngang: phục vụ − 2.320, 0,03 cm mỗi lượt; trục dọc: trễ, 0,0125 cm mỗi lượt.
\begin{scope}[shift={(10.6,0)}]
\draw[->] (0,0)--(3.95,0);
\draw[->] (0,0)--(0,3.35);
\foreach \vx/\so in {0/2.320,1.2/2.360,2.4/2.400,3.6/2.440}
  \draw (\vx,0)--(\vx,-0.07) node[below,nhan]{\so};
\foreach \vy/\so in {0/0,1/80,2/160,3/240}
  \draw (0,\vy)--(-0.07,\vy) node[left,nhan]{\so};
\draw (0.54,0.95)--(2.01,1.65)--(3.42,2.54);
\node[tron] at (0.54,0.95) {}; \node[nhan,below right] at (0.54,0.95) {30};
\node[thoi] at (2.01,1.65) {}; \node[nhan,below right] at (2.06,1.62) {120};
\node[tron] at (3.42,2.54) {}; \node[nhan,above left] at (3.42,2.54) {U};
\node at (1.8,-0.72) {Lượt phục vụ};
\node[anchor=south] at (1.9,3.4) {(c) Tập kiểm tra chính};
\end{scope}
\end{tikzpicture}
\caption{Lượt trễ quá 60 giây so với lời hứa đầu (trục dọc) theo lượt phục vụ (trục ngang) ở các mức
trần, lời hứa đầu không dự báo. Mỗi khung là một tập với thang riêng: (a) tập tắc đường cố định, 1.728
yêu cầu; (b) tập chỉnh trần buổi sáng, 864 yêu cầu; (c) tập kiểm tra chính, 3.672 yêu cầu. Số cạnh mỗi
điểm là trần tính bằng giây, U là điều phối không trần. Đường nối các mức của cùng một tập theo thứ tự
trần tăng dần; các tập không nối với nhau. Ở (a), ghép tự do (1.629 lượt phục vụ, 274 lượt trễ) nằm ngoài
khung; ở (b) chỉ vẽ các mức có đủ hai số đo ở Bảng~\ref{bang:c4-sang}. Hình thoi đen là trần 120 giây
trên tập kiểm tra chính.}
\label{hinh:c4-tran}
\end{figure}
% EV (a): D09 (C-30 1.534/0; C-60 1.540/0; C-120 1.563/23; U 1.584/47; FREE 1.629/274)
% EV (b): D01; D02 (C-30 660/47; C-90 664/50; C-120 666/51; U 665/50)
% EV (c): C01; C06; N-06 (C-30 2.338/76; C-120 2.387/132; U 2.434/203)
% Tọa độ: (a) x = (phục vụ − 1530) × 0,06, y = trễ × 0,05; (b) x = (phục vụ − 658) × 0,36, y = trễ × 0,05;
% (c) x = (phục vụ − 2320) × 0,03, y = trễ × 0,0125 (tính tay; xem NOTES §5).

Trên cả ba tập (Hình~\ref{hinh:c4-tran}), nới trần từ 30 lên 120 giây thêm lượt phục vụ và thêm lượt trễ
quá một phút. Số lượt không trễ quá một phút thì thêm 6 lượt trên tập tắc đường cố định, thêm 2 trên tập
buổi sáng và bớt 7 trên tập chính; ở ngưỡng hai phút, trần 120 giây không kém trần 30 giây ở cả ba tập.
Không mức trần nào tốt hơn ở mọi tiêu chí. Vì mỗi kịch bản chỉ chạy một lần ở mỗi mức và mức 120 giây được
chọn sau khi xem các tập phát triển, trần 120 giây trên tập chính là một điểm đo thêm trên đường đánh
đổi, không phải kết quả của một phép dò tìm.
% EV: D10; D03; L1a-09; D09; D02; L1a-02; K36; D11; B11; I06; K42

% -----------------------------------------------------------------------------------------------------
\section{Phần cộng dồn của luật C: kiểm toàn bộ bản ghi chặn}
\label{muc:c4-congdon}

Trong 98 lần chạy cũ có luật C dùng hệ số tắc đường đo, cả 1.256 bản ghi chặn do trần đều có một lần
dời đơn lẻ đã vượt trần của cấu hình tương ứng. Trên 34 lần chạy trần 120 giây của tập kiểm tra chính, 3
trong 524 bản ghi chỉ vượt trần khi cộng thêm phần trần đã dùng trước đó.
% EV: E01; K09; N-08; L1a-06

Một bản ghi chặn là \emph{chỉ do cộng dồn} khi lần dời mới không vượt $B$ nhưng phần đã dùng $L_r$ cộng
lần dời mới thì vượt; đó là dấu vết trực tiếp của việc luật C cộng các lần dời qua cả chuyến
(mục~\ref{muc:luatc}). Bảng~\ref{bang:c4-congdon} đếm các bản ghi này trên mọi bản ghi chặn của các lần
chạy đã kiểm.
% EV: A06; E05 (ví dụ 51,755 + 80,304 = 132,059 s > 120 s); E08; N-08 (quy tắc riêng + cumulative_decisive_check)

\begin{table}[htbp]
\centering
\caption{Kiểm toàn bộ bản ghi chặn do trần trên các lần chạy có luật C dùng hệ số tắc đường đo. Một bản ghi
ứng với một phương án bị loại, không phải một khách.}
\label{bang:c4-congdon}
\small
\begin{tabular}{@{}p{6.0cm}rrrr@{}}
\toprule
\textbf{Lần chạy} & \textbf{Số lần chạy} & \textbf{Bản ghi chặn} & \textbf{Chỉ do cộng dồn} &
\textbf{\begin{tabular}[b]{@{}r@{}}Khách phục vụ bị\\ dời từ hai lần\end{tabular}}\\
\midrule
Tập kiểm tra chính, trần 30 giây & 34 & 845 & 0 & 0\\
Tập chỉnh trần buổi sáng, trần 30--300 giây & 64 & 411 & 0 & 0\\
% EV: E01 (98 = 34 + 64; 1.256 = 845 + 411; 0 chỉ do cộng dồn); E02 (repeated = 0 trên 98 lần chạy)
Tập kiểm tra chính, trần 120 giây & 34 & 524 & 3 & 0\\
% EV: N-08; L1a-06; N-07 (0 khách ≥ 2 lần)
Chiều tuyết 15/11/2018, trần 30 giây & 8 & 163 & 0 & 0\\
% EV: N-26 (163 bản ghi, 0 chỉ do cộng dồn, ≥ 2 lần 0)
\midrule
Cộng & 140 & 1.943 & 3 & 0\\
% EV: N-15 (132 / 1.780 / 3; kể cả chiều tuyết 140 / 1.943 / 3); ADDENDUM-1 mục 0.3
\bottomrule
\end{tabular}
\end{table}

Cách đọc bảng: trên 34 lần chạy trần 120 giây có 524 phương án bị loại kèm bản ghi chặn do trần, 3 bản
ghi trong số đó chỉ do cộng dồn, và không khách được phục vụ nào bị dời từ hai lần.
% EV: N-08; N-07

Phạm vi kiểm là 132 lần chạy và 1.780 bản ghi (98 lần chạy cũ với 64.371 lượt quyết định được xác nhận,
và 34 lần chạy trần 120 giây), hay 140 lần chạy và 1.943 bản ghi nếu kể cả chiều tuyết; ``vượt trần'' là
vượt trần của chính cấu hình. 15 lần chạy hiệu chỉnh và 16 lần chạy đưa dự báo vào bộ điều phối chưa được
kiểm, nên kết quả không áp cho mọi lần chạy.
% EV: N-15 (SCOPE); ADDENDUM-1 mục 0.3 (15 lần chạy hiệu chỉnh, 16 thư mục C-30F không đọc); E01 (64.371); K09

Ba bản ghi chỉ do cộng dồn nằm ở ba kịch bản khung 8 giờ (ngày 14, 15 và 18/11), trong 3 lượt quyết định.
Ở bản ghi ngày 14/11, khách đã nhận đã dùng 25,844 giây của trần, và phương án mới sẽ dời giờ trả của
khách thêm 99,027 giây: mỗi phần dưới 120 giây, nhưng tổng 124,871 giây vượt trần. Khi không trần, chính
khách này bị dời hai lần đúng bằng hai lượng ấy. Trong 3 yêu cầu mới liên quan, 2 cuối cùng không được
nhận dưới trần 120 giây (một trong hai cũng không được nhận dưới trần 30 giây) và 1 được phục vụ. Vì bản
ghi dừng ở lỗi gặp trước và quỹ đạo các cấu hình khác nhau, điều này không chứng minh phần cộng dồn làm
mất hai yêu cầu đó.
% EV: L1a-06 (ba kịch bản w08 ngày 14, 15, 18/11); N-08 (3 quyết định; 2 không được nhận, 1 phục vụ); N-27 (25,844 + 99,027 = 124,871 s; 105,508 + 99,301 = 204,809 s); E07; A23; I08

Ngoài hệ số đo còn hai chỗ có dấu vết. Trên tập tắc đường cố định, nơi bộ kiểm ghi đủ mọi bản ghi thay
vì dừng ở lỗi gặp trước, có 2 bản ghi chỉ do cộng dồn ở trần 60 giây và 4 ở trần 120 giây; mọi yêu cầu
mới liên quan vẫn được phục vụ. Ở một kiểu tắc đường tự tạo có 4 bản ghi ở trần 120 giây, với 1 yêu cầu
liên quan bị từ chối.
% EV: E03; E04; A23 (NEW:13: tập tắc đường cố định bật chế độ ghi đủ; R23 #41)

Theo định nghĩa, vế cộng dồn vẫn là một nghĩa vụ khác: luật ``mỗi lần dời không quá 30 giây'' cho phép
ba lần dời, tổng 90 giây, còn luật C với trần 30 giây thì không. Khi không trần, 12 trong 2.434 khách bị
dời từ hai lần (Bảng~\ref{bang:c4-doidieuphoi}). Chấm lại các quỹ đạo không trần đã ghi, hai cách tính
chỉ khác nhau ở 4 trong 2.434 khách tại mức 120 giây và ở 0 khách tại 30, 60 và 300 giây. Phép sàng lọc
này không cho biết quỹ đạo của một luật khác sẽ ra sao. Luật mỗi lần một trần chưa được cài đặt và chưa
chạy, nên chưa biết bỏ vế cộng dồn có làm đổi quỹ đạo hay kết quả hay không.
% EV: E08; C20; E06 (0/0/4/0 trên 2.434 khách U); K42; E09; K08

Các mục trên so luật C với điều phối không trần khi giữ nguyên bộ chọn ưu tiên ổn định;
mục~\ref{muc:c4-vaitro} xét kết quả đó có phụ thuộc vào thứ tự tiêu chí của bộ chọn hay không.
% EV: A08; H01; H02
